\documentclass{article}

\usepackage[utf8]{inputenc}
\usepackage[T1]{fontenc}
\usepackage{amsmath,amssymb,amsfonts}
\usepackage{bm}
\usepackage[margin=1in]{geometry}
\usepackage{enumitem}

\newcommand{\vo}{\bm{o}}
\newcommand{\vz}{\bm{z}}
\newcommand{\va}{\bm{a}}
\newcommand{\vb}{\bm{b}}
\newcommand{\loss}{\mathcal{L}}
\newcommand{\E}{\mathbb{E}}

\title{Explanation: Why the AC-UL Prior Forces Physics-Grounded Encoding}
\author{}
\date{}

\begin{document}
\maketitle

\section{Setup}

We have two encoders:
\begin{itemize}[leftmargin=*]
  \item \textbf{Observation encoder:} $\vz_{k,\text{clean}} = E_\theta(\vo_k)$ --- maps each observation to a latent.
  \item \textbf{Action encoder:} $\vb_k = A_\theta(\va_k)$ --- maps each raw action to an action latent.
\end{itemize}

The dynamics prior must jointly \emph{denoise} the current action latent $\vb_t$ and the resulting future observation latent $\vz_{t+1,\text{clean}}$ from the latent history and language instruction:
\[
  \mathcal{H}_t^z = (\bm{l},\; \vz_{t-H+1,\text{clean}},\; \vb_{t-H},\; \vz_{t-H+2,\text{clean}},\; \vb_{t-H+1},\; \ldots,\; \vb_{t-1},\; \vz_{t,\text{clean}}).
\]
Here $\bm{l}$ is the language instruction (e.g., ``pick up the red block''), $\vz_{k,\text{clean}}$ are past observation latents, and $\vb_k$ are past action latents.

\textbf{Key autoregressive structure:} At the previous time step, the model denoised $(\vb_{t-1}, \vz_t)$---the action latent $\vb_{t-1}$ corresponds to the action that \emph{caused} the transition to observation $\vo_t$ (encoded as $\vz_t$). Now at time $t$, the model denoises $(\vb_t, \vz_{t+1})$: the action latent $\vb_t$ to execute now (decoded via $D_\theta^a(\vb_t)$ to get the raw action), and the frame $\vz_{t+1}$ that will \emph{result from} executing that action.

\textbf{Both modalities use the same diffusion objective.}
Given a shared noise level $\tau$, we noise both targets:
\begin{align}
  \vz_{t+1,\tau} &= \alpha_z(\tau)\,\vz_{t+1,\text{clean}} + \sigma_z(\tau)\,\epsilon_z, \\
  \vb_{t,\tau} &= \alpha_b(\tau)\,\vb_t + \sigma_b(\tau)\,\epsilon_b.
\end{align}

The joint prior denoises both simultaneously:
\[
  (\hat{\vb}_t, \hat{\vz}_{t+1}) = \hat{f}_\theta^{\text{dyn}}(\vb_{t,\tau}, \vz_{t+1,\tau}, \mathcal{H}_t^z).
\]

The joint prior loss is:
\begin{equation}
  \loss_z^{\text{joint}} = \underbrace{\|\vz_{t+1,\text{clean}} - \hat{\vz}_\theta\|^2}_{\text{obs.\ latent denoising (forward dynamics)}} + \gamma \cdot \underbrace{\|\vb_t - \hat{\vb}_\theta\|^2}_{\text{action latent denoising (inverse dynamics)}}.
  \label{eq:joint_prior}
\end{equation}

(The diffusion weighting $-\frac{d\lambda(\tau)}{d\tau}\frac{\exp\lambda(\tau)}{2}$ is omitted for clarity; it only rescales the MSE across noise levels.)

The two diffusion MSE terms play complementary roles:
\begin{itemize}[leftmargin=*]
  \item \textbf{Observation latent denoising (forward dynamics):} The prior must denoise $\vz_{t+1}$---what will happen after executing the action, given the history and language goal. This penalizes encoding anything unpredictable.
  \item \textbf{Action latent denoising (inverse dynamics):} The prior must denoise $\vb_t$---what action caused the transition $\vz_t \to \vz_{t+1}$. This prevents the observation encoder from collapsing to a trivially predictable but uninformative representation, and forces the action encoder to encode actions in a denoising-compatible form.
\end{itemize}

\section{Two Types of Information}

Consider what the observation encoder could put into $\vz_{t+1,\text{clean}}$.
There are two kinds of information:

\begin{enumerate}[leftmargin=*]
  \item \textbf{Predictable content} (e.g., gripper moved to grasp the red block because the language instruction said ``pick up the red block'' and the history shows the gripper approaching):
  The prior \emph{can} predict this from the latent history $\mathcal{H}_t^z$ (past observation latents, past action latents, and language goal), so the observation latent MSE stays low.
  This information is ``free''---it does not increase the prior loss.

  \item \textbf{Unpredictable content} (e.g., a shadow flickered, a background texture changed, camera noise):
  The prior \emph{cannot} predict this from $(\bm{l}, \vz_{t-H+1:t,\text{clean}}, \vb_{t-H:t-1})$ because nothing in the latent history or language instruction explains it.
  This spikes the observation latent MSE, increasing the prior loss.
\end{enumerate}

\section{The Trade-Off (Four-Way Pressure)}

The total Stage~1 loss is:
\begin{equation}
  \loss_{\text{S1}} = \underbrace{\loss_z^{\text{latent}}}_{\text{obs.\ latent diffusion}} + \gamma \cdot \underbrace{\loss_b}_{\text{action latent diffusion}} + c_{\text{lf}} \cdot \underbrace{\loss_x}_{\text{reconstruction}} + \underbrace{\loss_a^{\text{recon}}}_{\text{action reconstruction}}.
  \label{eq:total}
\end{equation}

\begin{itemize}[leftmargin=*]
  \item The \textbf{decoder loss} $\loss_x$ wants \emph{more} information in $\vz_{t+1,\text{clean}}$ (better reconstruction of $\vo_{t+1}$).
  \item The \textbf{observation latent prior loss} $\loss_z^{\text{latent}}$ penalizes information that the prior cannot denoise from the latent history and language.
  \item The \textbf{action latent prior loss} $\loss_b$ uses the same diffusion objective to ensure the action latent can be denoised from the latent history---requiring the latent to encode the causal effect of the action that \emph{caused} $\vo_{t+1}$.
  \item The \textbf{action reconstruction loss} $\loss_a^{\text{recon}} = \|\va_t - D_\theta^a(\vb_t)\|^2$ ensures the action encoder-decoder pair faithfully represents the action space.
  \item At equilibrium, both encoders keep only what the prior can denoise (physics that the interaction history and language goal explain), plus enough action-relevant detail to recover the causal action.
\end{itemize}

The information decomposition makes this precise:
\begin{equation}
  I(\vz_{t+1,\text{clean}};\, \vo_{t+1}) = \underbrace{I(\vz_{t+1,\text{clean}};\, \vo_{t+1} \mid \mathcal{H}_t^z)}_{\text{unpredictable (penalized by } \loss_z^{\text{latent}}\text{)}} + \underbrace{I(\vz_{t+1,\text{clean}};\, \mathcal{H}_t^z)}_{\text{predictable from latent history (free)}}.
  \label{eq:info_decomp}
\end{equation}

The observation latent prior loss upper-bounds the first term (see explanation2.pdf for the full proof).
The decoder loss incentivizes maximizing the total.
The action latent prior loss ensures the latent captures the causal effect of the action via the same diffusion objective.
Together, they push both encoders to maximize predictable, action-relevant information and discard unpredictable information.

\section{Why the Action Latent Diffusion Loss is Necessary}

\textbf{Without the action latent loss ($\gamma = 0$)}, the observation encoder could trivially minimize $\loss_z^{\text{latent}}$ by encoding a constant---perfectly predictable but completely uninformative.
The decoder loss would fight this, but the result would be a poor equilibrium where the latent captures some visual content without any guarantee of action-relevance.

\textbf{With the action latent diffusion loss ($\gamma > 0$)}, the prior must also denoise $\vb_t$ from $(\vb_{t,\tau}, \vz_{t+1,\tau}, \mathcal{H}_t^z)$.
This is an \emph{inverse dynamics} signal via diffusion: given the noised action latent, the noised observation latent, and the history, denoise the action that caused the transition.
The latent transition $\vz_t \to \vz_{t+1}$ must be informative enough that the action can be denoised from its noisy version.

Concretely, if $\va_t$ commands the gripper to close on the red block (as instructed by $\bm{l}$), the action encoder must encode this into $\vb_t$, and the observation latent $\vz_{t+1,\text{clean}}$ must encode the resulting contact geometry such that the prior can jointly denoise both.
The denoised action latent $\hat{\vb}_t$ is then decoded via $D_\theta^a(\hat{\vb}_t)$ to recover the raw action $\hat{\va}_t$.

The rate-distortion equilibrium is:
\begin{equation}
  \vz_{t+1,\text{clean}}^* = \arg\min_{\vz} \Big[
    \underbrace{I(\vz;\, \vo_{t+1} \mid \mathcal{H}_t^z)}_{\substack{\text{rate: unpredictable bits} \\ \text{(minimize)}}}
    + c_{\text{lf}} \cdot \underbrace{D(\vo_{t+1},\, \hat{\vo}_\theta)}_{\substack{\text{distortion: reconstruction} \\ \text{(minimize)}}}
    - \gamma \cdot \underbrace{I(\vb_t;\, \vz \mid \mathcal{H}_t^z)}_{\substack{\text{action-obs coupling} \\ \text{(maximize)}}}
  \Big].
  \label{eq:rd_action}
\end{equation}

The three terms create a balanced pressure:
\begin{enumerate}[leftmargin=*]
  \item \textbf{Rate} ($\loss_z^{\text{latent}}$): Compress out everything unpredictable from latent history and language.
  \item \textbf{Distortion} ($\loss_x$): Retain enough information for good reconstruction.
  \item \textbf{Action-obs coupling} ($\loss_b$): Ensure the observation latent and action latent are mutually informative---the latent transition encodes the causal effect of the action, and the action latent is jointly denoised.
\end{enumerate}

The result is a latent that encodes \emph{exactly the action-predictable physical geometry}:
gripper poses, object displacements, contact states---while discarding textures, lighting, and background clutter that neither the action history nor the language instruction can explain.

\section{Caveat}

The claim is somewhat idealized. In practice:
\begin{itemize}[leftmargin=*]
  \item The prior has finite capacity, so some predictable information may still be penalized if the prior network cannot learn the dynamics well enough.
  \item The loss factor $c_{\text{lf}}$ controls the trade-off: large $c_{\text{lf}}$ lets more information through (decoder wins), small $c_{\text{lf}}$ compresses harder (prior wins).
  \item ``Mathematically forced'' is strong language---it is more accurate to say the loss \emph{incentivizes} the encoders in this direction.
    The bound is tight only when $w(\lambda) = 1$, and the loss factor loosens it further.
\end{itemize}

\end{document}
